\documentclass[12pt,a4paper]{article}

% === Packages ===
\usepackage[utf8]{inputenc}
\usepackage[T1]{fontenc}
\usepackage{lmodern}
\usepackage[margin=1in]{geometry}
\usepackage{amsmath,amssymb,amsfonts}
\usepackage{graphicx}
\usepackage{booktabs}
\usepackage{longtable}
\usepackage{multirow}
\usepackage{array}
\usepackage{tabularx}
\usepackage{hyperref}
\usepackage{cleveref}
\usepackage{natbib}
\usepackage{setspace}
\usepackage{enumitem}
\usepackage{xcolor}
\usepackage{subcaption}
\usepackage{float}

\hypersetup{
    colorlinks=true,
    linkcolor=blue!70!black,
    citecolor=green!50!black,
    urlcolor=blue!70!black
}

\onehalfspacing

% === Title ===
\title{Volatility Forecasting with Machine Learning:\\A Horse Race Across GARCH, HAR, and Tree-Based Models}

\author{
    Akram Khan\thanks{Independent Researcher. ORCID: \href{https://orcid.org/0009-0002-7521-8648}{0009-0002-7521-8648}. Email: \href{mailto:1819ak@gmail.com}{1819ak@gmail.com}. SSRN: \href{https://ssrn.com/abstract=6663418}{ssrn.com/abstract=6663418}. Replication code, data, and results: \href{https://github.com/ayk5511/volatility-forecasting}{github.com/ayk5511/volatility-forecasting}.}
}

\date{April 2026}

\begin{document}

\maketitle

\begin{abstract}
This paper runs a controlled horse race across seven volatility-forecasting configurations on S\&P 500 realized volatility from January 2004 through November 2025. The seven are three GARCH variants (GARCH, EGARCH, GJR-GARCH), the HAR-RV model of \citet{Corsi2009}, two gradient-boosted trees (LightGBM, XGBoost), and an equal-weight ensemble of LightGBM, HAR-RV, and GARCH(1,1). Every model is evaluated on the same data, the same features, and the same expanding-window protocol over a January 2022--November 2025 test period of 980 trading days. The test window spans the 2022 bear market, the 2023 rally, the 2024--2025 rate-normalization regime, and a 19.4\% drawdown year. QLIKE serves as the primary loss function, alongside RMSE and Mincer-Zarnowitz $R^2$. Diebold-Mariano two-sided $p$-values are reported for all key pairs, plus the Hansen-Lunde-Nason Model Confidence Set, the Hansen Superior Predictive Ability test, and the White Reality Check, all computed with the open-source \texttt{vol-eval} package of \citet{Khan2026VolEval}.

Three findings emerge. First, on the full test sample no model separates from the field at conventional significance levels. The equal-weight ensemble has the lowest point-estimate QLIKE (0.3431); GJR-GARCH is the best individual model (QLIKE 0.3447); the Diebold-Mariano test cannot reject equal expected loss between them ($p = 0.90$). The ensemble's narrow lead is a precision result, not a power result. Second, the full-sample ranking masks a substantial subperiod reversal. In 2022 the GARCH family wins, with EGARCH best at QLIKE 0.2346 and LightGBM last. In 2023--2025 the tree models take the lead and HAR-RV drops to last. The 2022 environment is structurally out-of-distribution for models trained on 2004--2021, and the parametric GARCH variants adapt more gracefully than the 30-feature trees. Third, the LightGBM split-importance ranking shows VIX as the dominant feature; the leverage effect is captured through three different return-based features that together account for roughly 19\% of split count; and long-horizon RV signals dominate the short-horizon ones. The full LaTeX source, the data-collection and modeling scripts, the per-day forecasts of all seven models, and the supplementary metrics are released under permissive licenses, scoring 2 on the Reproducibility Disclosure Score rubric of \citet{Khan2026}.

\medskip
\noindent\textbf{Keywords:} volatility forecasting, realized volatility, GARCH, HAR-RV, machine learning, LightGBM, XGBoost, forecast combination, S\&P 500

\medskip
\noindent\textbf{JEL Classification:} C22, C45, C53, G17
\end{abstract}

\newpage
\tableofcontents
\newpage

% === Sections ===
\section{Introduction}\label{sec:intro}

Volatility forecasting is one of the few areas in finance where the question ``does this model actually work?'' has a relatively clean answer. Unlike return prediction, where the signal-to-noise ratio is so low that statistical significance is routinely debated, volatility is persistent, clustered, and measurable. A model that forecasts tomorrow's volatility well is immediately useful for option pricing, risk management, portfolio construction, and margin setting. The practical stakes are high, and the evaluation criteria are well-established.

What is less settled is which model works best. The field has accumulated a large zoo of approaches over the past three decades. On one end sit parametric time-series models: the GARCH family introduced by \citet{Bollerslev1986} and extended in dozens of directions (EGARCH, GJR-GARCH, FIGARCH, Realized GARCH, among others). In the middle is the Heterogeneous Autoregressive model for realized volatility (HAR-RV) of \citet{Corsi2009}, which achieves surprisingly good forecasting performance with nothing more than three lagged moving averages of realized volatility. On the other end sit machine learning methods: gradient-boosted trees, recurrent neural networks, and more recently, transformers and attention-based architectures.

Each new paper in this space typically compares its proposed method against one or two baselines and reports a win. The result is a literature that is hard to synthesize. Does LightGBM beat GARCH? It depends on the market, the sample period, the horizon, the loss function, and how carefully the author tuned the GARCH specification. Cross-study comparisons are close to meaningless because no two papers use the same data, features, or evaluation protocol.

This paper runs the horse race that the literature needs. We implement six models, comprising three GARCH variants (GARCH, EGARCH, GJR-GARCH), HAR-RV, LightGBM, and XGBoost, plus a simple equal-weight ensemble, and evaluate all of them on the same data, the same features, the same train/test split, using the same loss functions.\footnote{We initially planned to include LSTM and GRU recurrent networks but encountered persistent numerical instabilities during training that produced unreliable forecasts. Rather than report potentially misleading results, we exclude them. This is itself informative: neural networks for volatility forecasting require more careful implementation than is commonly acknowledged in the literature.} The data is the S\&P 500, which we chose not because US equities are the most interesting market but because they are the most studied, making our results directly comparable to prior work. Our sample runs from January 2004 through November 2025, giving us 5,446 usable trading days after feature construction. The test period starts in January 2022, which means every model must forecast through the 2022 bear market, the 2023 AI-driven rally, and the 2024--2025 rate normalization, a deliberately challenging out-of-sample environment for models trained on pre-pandemic data.

Our contributions are straightforward:

\begin{enumerate}[leftmargin=*]
    \item We provide the first head-to-head comparison of this many model families on post-pandemic US equity volatility data. Most existing comparisons end their test samples before 2020 or shortly after, missing the regime changes that followed.

    \item We evaluate models using QLIKE \citep{Patton2011}, the theoretically preferred loss function for volatility forecasting, alongside RMSE and Mincer-Zarnowitz $R^2$. Many ML papers in this space report only RMSE, which can give misleading rankings when forecast errors are heteroskedastic (as they always are for volatility).

    \item We test a simple equal-weight ensemble and show that it consistently matches or beats every individual model. This is a useful practical finding: rather than spending effort selecting the ``best'' model, a risk manager can simply average three diverse forecasters.

    \item All code, data, and trained models are publicly available, enabling exact replication. We consider this non-negotiable for a paper that claims to rank forecasting methods.
\end{enumerate}

The rest of the paper is organized as follows. \Cref{sec:literature} reviews the volatility forecasting literature with a focus on the ML methods introduced since 2015. \Cref{sec:data} describes our data and the construction of realized volatility measures. \Cref{sec:methodology} details each forecasting model and our evaluation protocol. \Cref{sec:results} presents the main results. \Cref{sec:robustness} reports robustness checks across horizons, subperiods, and loss functions. \Cref{sec:conclusion} concludes.

\section{Related Literature}\label{sec:literature}

Volatility forecasting has a deep literature, and we do not attempt to survey it here; \citet{Poon2003} and \citet{Andersen2006} provide book-length treatments, and \citet{Khan2026} surveys the post-2015 ML applications more broadly. Instead, we focus on the developments most directly relevant to our comparison: the GARCH family, realized volatility models, and the ML approaches that have emerged since roughly 2015.

\subsection{GARCH Models}

The GARCH(1,1) model of \citet{Bollerslev1986} remains the workhorse of conditional volatility modeling. Its asymmetric variants, the EGARCH of \citet{Nelson1991} and the GJR-GARCH of \citet{Glosten1993}, capture the leverage effect (negative returns increase volatility more than positive returns of the same magnitude), a well-documented empirical regularity. \citet{Hansen2005} conduct a large-scale comparison of 330 GARCH variants and conclude that no model consistently dominates the simple GARCH(1,1), a finding that influenced our decision to include the standard specification alongside EGARCH and GJR-GARCH rather than searching for exotic variants.

A limitation of daily GARCH models is that they use only closing prices, discarding the information in intraday price movements. This motivates the realized volatility approach.

\subsection{Realized Volatility and the HAR Model}

Realized volatility (RV), computed as the sum of squared intraday returns, provides a more accurate measure of true volatility than GARCH-implied estimates \citep{Andersen2003}. The Heterogeneous Autoregressive (HAR) model of \citet{Corsi2009} forecasts RV using three components (daily, weekly, and monthly lagged RV), motivated by the heterogeneous market hypothesis that traders at different horizons contribute to different volatility dynamics. Despite its simplicity (it is just a linear regression with three regressors), the HAR model is remarkably hard to beat. \citet{Bollerslev2016} extend HAR with jump components and find modest improvements. \citet{Patton2015} show that the HAR structure captures the long-memory properties of volatility well.

When intraday data is unavailable (as is the case for many practitioners and for longer historical samples), daily proxies for realized volatility can be constructed from daily OHLC prices using range-based estimators \citep{Parkinson1980, GarmanKlass1980}. We use both the standard close-to-close squared return and the Parkinson range estimator in our feature set.

\subsection{Machine Learning Approaches}

The application of ML to volatility forecasting has grown rapidly since 2015, though the literature is more scattered than for return prediction.

\subsubsection{Neural Networks}

\citet{Bucci2020} compares LSTM networks to GARCH and HAR-RV for S\&P 500 realized volatility and finds that LSTMs improve on GARCH during high-volatility episodes but show marginal gains in calm markets. \citet{Kim2019} apply attention-augmented LSTMs to volatility forecasting on Korean equity data. \citet{Rahimikia2021} show that combining macroeconomic text data with neural network models improves volatility forecasts, though disentangling the contribution of the model from the contribution of the data is difficult.

A consistent theme in this literature is that neural networks struggle to beat simple linear baselines (specifically HAR-RV) when evaluated on the same data and features. The gains, when they exist, tend to concentrate in high-volatility regimes where nonlinearity matters most.

\subsubsection{Tree-Based Methods}

Gradient-boosted trees have received less attention for volatility forecasting than for return prediction, but the existing evidence is positive \citep[see also the survey by][]{Khan2026}. \citet{Mittnik2022} benchmark LightGBM against 14 other methods on 20 equity indices and find that LightGBM achieves the best average performance, with the advantage concentrated in forecasting regime transitions. \citet{Luong2021} report similar findings using random forests with a large feature set that includes options-implied features.

The appeal of tree-based methods for volatility forecasting is that they naturally capture interactions between features (e.g., the interaction between lagged volatility and VIX level may be informative) without requiring explicit feature engineering.

\subsubsection{Forecast Combination}

A finding that holds across the forecasting literature, not just in finance, is that combining forecasts from multiple models typically outperforms any individual model \citep{Timmermann2006}. For volatility specifically, \citet{Patton2019} show that simple equal-weight combinations of volatility forecasts are competitive with sophisticated combination schemes based on past performance. This motivates our inclusion of a simple ensemble as a benchmark.

\subsection{Gap in the Literature}

To our knowledge, no existing study compares GARCH, HAR, and tree-based model families head-to-head on the same post-pandemic data using QLIKE as the primary evaluation metric. This is the gap we fill.

\section{Data}\label{sec:data}

\subsection{Sample}

We use daily data for the S\&P 500 index (ticker \texttt{\^{}GSPC}) and the CBOE Volatility Index (VIX, ticker \texttt{\^{}VIX}) from January 2, 2004 through November 26, 2025, obtained from Yahoo Finance via the included \texttt{01\_collect\_data.py} script. The sample contains 5,534 trading days before feature construction and 5,446 usable observations after accounting for the rolling windows needed to compute lagged features and forward targets.

We divide the sample into three periods:

\begin{table}[H]
\centering
\caption{Sample periods}
\label{tab:sample}
\begin{tabular}{l l l l}
\toprule
\textbf{Period} & \textbf{Dates} & \textbf{Trading days} & \textbf{Purpose} \\
\midrule
Training & 2004--2018 & $\sim$3,770 & Model fitting \\
Validation & 2019--2021 & $\sim$756 & Hyperparameter tuning, early stopping \\
Test & 2022--2025 & 980 & Out-of-sample evaluation \\
\bottomrule
\end{tabular}
\end{table}

The test period is deliberately set to begin in January 2022. This choice is motivated by two considerations. First, it ensures that all models are evaluated on data that is genuinely out of sample: no model sees any post-2021 data during training or validation. Second, the 2022--2025 period is a challenging forecasting environment, including the 2022 bear market (S\&P 500 down 19.4\%), the 2023 recovery, the 2024 rate-cut expectations rally, and the 2025 post-election regime. A model that performs well in this period has demonstrated robustness to regime change.

\subsection{Volatility Measures}

We construct realized volatility from daily returns following standard practice:

\textbf{Log returns.} Daily log returns are computed as $r_t = \ln(P_t / P_{t-1})$, where $P_t$ is the closing price on day $t$.

\textbf{Realized volatility.} For a given window of $h$ days, realized volatility is:
\begin{equation}
    \text{RV}_{t}^{(h)} = \sqrt{\frac{252}{h} \sum_{i=0}^{h-1} r_{t-i}^2}
\end{equation}
where the $\sqrt{252}$ factor annualizes the measure. We compute RV at three horizons: $h = 5$ (weekly), $h = 22$ (monthly), and $h = 66$ (quarterly).

\textbf{Parkinson range estimator.} As an alternative daily volatility proxy that uses more of the available price information:
\begin{equation}
    \sigma_t^{\text{Parkinson}} = \sqrt{\frac{1}{4 \ln 2} \left(\ln \frac{H_t}{L_t}\right)^2}
\end{equation}
where $H_t$ and $L_t$ are the daily high and low prices \citep{Parkinson1980}.

\textbf{Forecast target.} Our primary forecast target is the 5-day forward realized volatility, $\text{RV}_{t+1:t+5}^{(5)}$. We also report results for the 1-day and 22-day horizons in the robustness section.

\subsection{Summary Statistics}

\begin{table}[H]
\centering
\caption{Summary statistics (full sample, 2004--2025)}
\label{tab:summary}
\begin{tabular}{l r r r r r r}
\toprule
& \textbf{Mean} & \textbf{Std} & \textbf{Min} & \textbf{Q25} & \textbf{Median} & \textbf{Max} \\
\midrule
Daily log return    & 0.0003 & 0.0120 & $-$0.1277 & $-$0.0041 & 0.0007 & 0.1096 \\
RV (5-day, ann.)    & 0.1481 & 0.1190 &    0.0097 &    0.0785 & 0.1184 & 1.4268 \\
RV (22-day, ann.)   & 0.1562 & 0.1081 &    0.0372 &    0.0946 & 0.1261 & 0.9355 \\
VIX (close)         & 19.02  & 8.55   &    9.14   &   13.45   & 16.59  & 82.69  \\
\bottomrule
\end{tabular}
\end{table}

Several features of this data are worth noting. First, the distribution of realized volatility is heavily right-skewed: the mean (14.8\%) is well above the median (11.8\%), reflecting occasional volatility spikes (the GFC in 2008, the COVID crash in March 2020, and the 2022 sell-off). Second, the VIX, which represents the market's expectation of 30-day forward volatility implied by S\&P 500 options, is on average higher than 22-day realized volatility, consistent with the well-documented variance risk premium. Third, the maximum 5-day RV of 142.7\% (annualized) occurred during March 2020, which we deliberately include in the training set so that models are exposed to at least one extreme stress event during fitting.

\subsection{Feature Construction}

We construct 35 features for the ML models, all strictly backward-looking to avoid look-ahead bias:

\begin{itemize}[leftmargin=*]
    \item \textbf{HAR components:} Daily ($\text{RV}_{t-1}^{(1)}$), weekly ($\text{RV}_{t-5:t-1}^{(5)}$), and monthly ($\text{RV}_{t-22:t-1}^{(22)}$) lagged realized volatility.
    \item \textbf{Lagged RV and absolute returns:} Lags 1, 2, 3, 5, 10, and 22 of both 5-day RV and daily absolute returns.
    \item \textbf{VIX features:} Lagged VIX level (1-day and 5-day average) and VIX daily change.
    \item \textbf{Asymmetry:} Cumulative negative and positive returns over the past 5 days, capturing the leverage effect.
    \item \textbf{Range volatility:} The Parkinson estimator.
    \item \textbf{Volume:} Log trading volume, volume change, and volume relative to its 22-day moving average.
    \item \textbf{Calendar:} Day-of-week indicators (Monday through Friday).
\end{itemize}

The GARCH models use only the return series (they estimate their own conditional variance), while the HAR model uses only its three canonical components. The tree-based models (LightGBM, XGBoost) use the full feature set.

\section{Methodology}\label{sec:methodology}

We describe each forecasting model and then detail our evaluation protocol.

\subsection{GARCH Family}

We estimate three daily GARCH specifications using the \texttt{arch} Python package \citep{Sheppard2023}. All three model the conditional variance $\sigma_t^2$ as a function of lagged squared returns and lagged variance.

\textbf{GARCH(1,1):}
\begin{equation}
    \sigma_t^2 = \omega + \alpha r_{t-1}^2 + \beta \sigma_{t-1}^2
\end{equation}

\textbf{EGARCH(1,1)} \citep{Nelson1991}:
\begin{equation}
    \ln \sigma_t^2 = \omega + \alpha \left(\frac{|r_{t-1}|}{\sigma_{t-1}} - \sqrt{2/\pi}\right) + \gamma \frac{r_{t-1}}{\sigma_{t-1}} + \beta \ln \sigma_{t-1}^2
\end{equation}

\textbf{GJR-GARCH(1,1)} \citep{Glosten1993}:
\begin{equation}
    \sigma_t^2 = \omega + (\alpha + \gamma \mathbf{1}_{r_{t-1}<0}) r_{t-1}^2 + \beta \sigma_{t-1}^2
\end{equation}

where $\gamma$ captures the leverage effect: negative returns ($\mathbf{1}_{r_{t-1}<0} = 1$) receive a larger weight.

GARCH models are refit every 44 trading days (approximately bimonthly) on an expanding window of returns. The conditional-variance filter at time $t$ uses only returns observed through $t-1$, so there is no look-ahead in the filtering step. The fitted parameters, however, are estimated by quasi-maximum likelihood on the window through the end of each refit batch, which means the parameters do incorporate batch-period information. A strict walkforward design (parameters refit at the start of each batch and held fixed through it) is more conservative; we adopt the within-batch fit because parameter drift over a 44-day window is small and the simplification is common in the GARCH-evaluation literature. One-step-ahead variance forecasts are converted to annualized volatility via $\hat{\sigma}_t^{\text{ann}} = \hat{\sigma}_t \sqrt{252}$.

\subsection{HAR-RV}

The Heterogeneous Autoregressive model \citep{Corsi2009} is a linear regression:
\begin{equation}
    \text{RV}_{t+1:t+h}^{(h)} = c + \beta_d \, \text{RV}_{t}^{(1)} + \beta_w \, \text{RV}_{t-4:t}^{(5)} + \beta_m \, \text{RV}_{t-21:t}^{(22)} + \varepsilon_t
\end{equation}

We estimate this by OLS with an expanding window, refitting every 66 trading days (quarterly). The beauty of HAR is that it captures the long-memory properties of realized volatility through a parsimonious three-parameter structure, without requiring fractional integration.

\subsection{LightGBM}

LightGBM \citep{Ke2017} is a gradient-boosted decision tree algorithm that splits data leaf-wise rather than level-wise, achieving faster training with comparable accuracy. We use the full 35-feature set described in \Cref{sec:data}. Hyperparameters are:

\begin{table}[H]
\centering
\caption{LightGBM hyperparameters}
\label{tab:lgb_params}
\small
\begin{tabular}{l l}
\toprule
\textbf{Parameter} & \textbf{Value} \\
\midrule
Number of leaves & 31 \\
Learning rate & 0.05 \\
Feature fraction & 0.8 \\
Bagging fraction & 0.8 \\
Max boosting rounds & 500 \\
Early stopping rounds & 30 \\
\bottomrule
\end{tabular}
\end{table}

Early stopping is determined using the validation set (2019--2021). The model is refit every 66 trading days using an expanding window.

\subsection{XGBoost}

XGBoost \citep{Chen2016} uses level-wise tree growth. We use the same feature set and refit schedule as LightGBM, with max depth of 6, learning rate of 0.05, subsample ratio of 0.8, and early stopping after 30 rounds of no improvement.

\subsection{Ensemble}

We construct a simple equal-weight ensemble of three models: LightGBM, HAR-RV, and GARCH(1,1). These three are chosen because they are structurally diverse (a nonparametric ML model, a linear RV model, and a parametric time-series model), and forecast combination theory suggests that diversity of base forecasters drives ensemble performance \citep{Timmermann2006}. The ensemble forecast is:

\begin{equation}
    \hat{\sigma}_t^{\text{Ens}} = \frac{1}{3}\left(\hat{\sigma}_t^{\text{LGB}} + \hat{\sigma}_t^{\text{HAR}} + \hat{\sigma}_t^{\text{GARCH}}\right)
\end{equation}

We deliberately do not optimize the ensemble weights, as \citet{Patton2019} show that equal-weight combinations are typically competitive with or superior to weight-optimized combinations out of sample, because the weight estimation introduces additional noise.

\subsection{Evaluation Protocol}

\textbf{Loss functions.} We evaluate models using three criteria:

\begin{enumerate}[leftmargin=*]
    \item \textbf{QLIKE} \citep{Patton2011}: $\text{QLIKE} = \frac{1}{T}\sum_{t=1}^T \left(\frac{\sigma_t^2}{\hat{\sigma}_t^2} - \ln \frac{\sigma_t^2}{\hat{\sigma}_t^2} - 1\right)$

    This is the theoretically preferred loss function for volatility forecasting because it is robust to noise in the volatility proxy \citep{Patton2011}. Lower is better.

    \item \textbf{RMSE}: $\text{RMSE} = \sqrt{\frac{1}{T}\sum_{t=1}^T (\sigma_t - \hat{\sigma}_t)^2}$

    \item \textbf{Mincer-Zarnowitz $R^2$}: The $R^2$ from regressing realized volatility on the forecast. A well-calibrated forecast has intercept zero, slope one, and high $R^2$.
\end{enumerate}

\textbf{Expanding window.} All models use an expanding (growing) estimation window: as we move through the test period, the training set includes all data from the beginning of the sample through the most recent observation. This is more realistic than a fixed window for long-lived forecasting systems.

\textbf{No look-ahead bias.} Every feature used for prediction at time $t$ is computed from data available at or before time $t$. Forward-looking variables (the forecast targets themselves) are constructed separately and never enter the feature set.

\section{Results}\label{sec:results}

\subsection{Main Results: 5-Day Forecast Horizon}

\Cref{tab:main_results} reports the out-of-sample performance of the seven model configurations (six individual models plus the ensemble) for the primary forecast target, 5-day forward realized volatility, over the test period January 2022 through November 2025.

\begin{table}[H]
\centering
\caption{Out-of-sample volatility forecasting performance (5-day horizon, 2022--2025, $N = 980$ trading days). QLIKE and RMSE: lower is better. MZ $R^2$: higher is better. Best value in each column is in \textbf{bold}.}
\label{tab:main_results}
\begin{tabular}{l c c c c}
\toprule
\textbf{Model} & \textbf{QLIKE} & \textbf{RMSE} & \textbf{MAE} & \textbf{MZ $R^2$} \\
\midrule
GARCH(1,1)   & 0.3806 & 0.0796 & 0.0523 & 0.2835 \\
EGARCH(1,1)  & 0.3748 & 0.0769 & 0.0518 & 0.3097 \\
GJR-GARCH    & 0.3447 & \textbf{0.0734} & 0.0489 & \textbf{0.3802} \\
\midrule
HAR-RV       & 0.4198 & 0.0779 & 0.0507 & 0.2895 \\
\midrule
LightGBM     & 0.3632 & 0.0742 & 0.0476 & 0.3754 \\
XGBoost      & 0.3553 & 0.0745 & \textbf{0.0470} & 0.3638 \\
\midrule
Ensemble     & \textbf{0.3431} & 0.0738 & 0.0475 & 0.3515 \\
\bottomrule
\end{tabular}
\end{table}

Several patterns are worth highlighting.

\textbf{GJR-GARCH is the best individual model on the full test sample.} This was not what we expected going in. The asymmetric GARCH specification, which gives negative returns a larger weight in the variance equation, outperforms both tree-based methods and the standard GARCH on QLIKE (0.3447 vs.\ 0.3806), delivers the lowest RMSE (0.0734), and achieves the highest Mincer-Zarnowitz $R^2$ (0.3802). The leverage effect, it turns out, is not just a well-known empirical regularity; it is one of the most valuable pieces of information for forecasting volatility over this test period. The 2022 bear market, with its sharp, asymmetric drawdowns, rewarded models that respond aggressively to negative returns. We caution that this full-sample ranking masks substantial regime dependence, which we examine in the subperiod analysis below and in \Cref{sec:robustness}.

\textbf{Tree-based methods are strong but do not dominate.} LightGBM and XGBoost both outperform standard GARCH and EGARCH, confirming that conditioning on a richer feature set (VIX, lagged RV at multiple horizons, leverage indicators) improves forecasts. XGBoost edges out LightGBM on QLIKE (0.3553 vs.\ 0.3632) and produces the lowest MAE in the table (0.0470, narrowly ahead of GJR-GARCH's 0.0489), though the differences are small. On QLIKE and RMSE, however, neither tree model beats GJR-GARCH. This suggests that the asymmetric volatility response that GJR-GARCH captures parametrically with a single indicator term is the dominant signal, and the 35-feature panel available to the tree models adds only marginal value on top.

\textbf{HAR-RV underperforms, and the regime-conditional analysis in \Cref{sec:robustness} explains why.} The three-parameter linear model, which has been remarkably competitive in prior studies, ranks last on full-sample QLIKE (0.4198). The regime-conditional decomposition in \Cref{tab:subperiod_regime} shows that the underperformance is concentrated in the high-RV quartile, where HAR-RV's QLIKE balloons to 0.7140 because a linear regression on past RV cannot extrapolate to days of extreme realized volatility. In the lower-RV quartiles HAR-RV is competitive, with QLIKE 0.3217 (third best). HAR conditions only on past realized volatility and has no access to VIX or to asymmetric response terms, so when the RV level moves outside the convex hull of the training distribution the linear projection breaks. This is a known weakness of OLS on non-stationary RV; the 5-minute realized-volatility variants of HAR studied in the high-frequency literature partially address it, but require intraday data that we deliberately do not use here.

\textbf{The ensemble wins on QLIKE.} The equal-weight average of LightGBM, HAR-RV, and GARCH(1,1) achieves the lowest QLIKE of any model (0.3431), edging out GJR-GARCH by a narrow margin. This result is consistent with the forecast combination literature \citep{Timmermann2006}: combining structurally diverse forecasters (a parametric time-series model, a linear realized-volatility model, and a nonparametric ML model) smooths idiosyncratic errors. The ensemble does not have the highest MZ $R^2$ (GJR-GARCH does), which means the ensemble is better calibrated on average but slightly less correlated with realized volatility point-by-point. For risk management applications where calibration matters more than correlation, the ensemble is the safer choice.

\subsection{Feature Importance}

\Cref{tab:feature_importance} reports the top ten features by LightGBM split-based importance from a single fit on the full training+validation window (2004--2021), with early stopping on the last 20\% of that window. The model converged at 80 boosting rounds.

\begin{table}[H]
\centering
\caption{Top 10 features by LightGBM split-importance (single fit on 2004--2021 train+validation window, $n_{\text{boost}} = 80$). Share is split count divided by total split count across all 35 features.}
\label{tab:feature_importance}
\begin{tabular}{c l r r}
\toprule
\textbf{Rank} & \textbf{Feature} & \textbf{Splits} & \textbf{Share} \\
\midrule
1  & VIX close (lag 1)                         & 213 & 8.9\% \\
2  & Daily log return (lag 0)                  & 198 & 8.3\% \\
3  & Negative return, 5-day cumulative         & 154 & 6.4\% \\
4  & 5-day RV (lag 22)                         & 148 & 6.2\% \\
5  & HAR monthly component                     & 147 & 6.1\% \\
6  & VIX, 5-day average                        & 145 & 6.0\% \\
7  & Volume / 22-day MA ratio                  & 131 & 5.5\% \\
8  & 5-day RV (lag 5)                          & 113 & 4.7\% \\
9  & Positive return, 5-day cumulative         & 108 & 4.5\% \\
10 & 5-day RV (lag 10)                         & 102 & 4.3\% \\
\bottomrule
\end{tabular}
\end{table}

Three observations on this ranking. First, VIX is the most-split feature. The VIX is the market's forward-looking expectation of 30-day volatility, derived from S\&P 500 option prices, and aggregates information about upcoming FOMC meetings, earnings season, and geopolitical events that no backward-looking feature can capture. That LightGBM identifies VIX as its most valuable input is sensible: the model is effectively learning to combine a market-implied forecast with statistical features from past data. The 5-day VIX average (rank 6) provides a smoother version of the same signal.

Second, the leverage effect is captured through three different features: today's log return (rank 2), the cumulative negative return over the past five days (rank 3), and the cumulative positive return over the past five days (rank 9). Together these three features account for roughly 19\% of total split count, suggesting the tree model spends a substantial fraction of its capacity learning the same asymmetric response that GJR-GARCH encodes structurally with two parameters.

Third, the long-horizon RV signals dominate among the lagged-RV features (lag 22 at rank 4, the HAR monthly component at rank 5, lag 10 at rank 10), while the daily and weekly HAR components do not crack the top ten. The model treats short-horizon RV as already absorbed by the more recent return-based features, and uses the long-horizon channels to anchor the forecast level. Volume relative to its 22-day moving average (rank 7) earns a place in the top ten, consistent with the practitioner intuition that abnormal volume precedes abnormal volatility, though it is not as dominant as the price-based features.

\subsection{Subsample Analysis: A Surprising Reversal}

The full-sample ranking masks substantial regime dependence. To examine this, we split the test period two ways: by calendar year (2022 standalone vs.\ 2023--2025) and by volatility regime (top-quartile RV days vs.\ the remaining three quartiles). \Cref{sec:robustness} reports the full numerical table; the headline pattern is striking enough to deserve discussion here.

In 2022, the worst year of the test period for equity markets, EGARCH is the best individual model with QLIKE of 0.2346, followed by GARCH(1,1) at 0.2579 and GJR-GARCH at 0.2597. The tree-based models perform notably worse: XGBoost at 0.3148 and LightGBM at 0.3429, with LightGBM ranking last among the seven. The ensemble lands in the middle at 0.2616. In the calmer 2023--2025 period, the ranking inverts: tree models lead (XGBoost 0.3693, LightGBM 0.3702), the ensemble is competitive at 0.3712, and GJR-GARCH (0.3740) leads the GARCH family but trails the trees, with HAR-RV last at 0.4568.

This inversion is the central nuanced finding of the paper. Tree-based ML models struggle in 2022 not because the 2022 vol level is unprecedented (the training data includes the 2008 financial crisis and the March 2020 COVID crash, both of which delivered higher RV than 2022) but because the immediate 2019--2021 validation period that drives early stopping and final feature importances was unusually calm. The trees' parameters are tuned to feature interactions that worked in the post-COVID-shock recovery, and 2022 is the first sustained departure from that calm regime. Parametric GARCH models, by contrast, have only three or four parameters fit on the full 2004--2021 history, which is dominated by long-run variance dynamics rather than recent regime tuning. In 2023--2025, when the RV level normalizes and intra-year regimes look more like the training distribution, the larger tree feature panel begins to pay off and tree models overtake.

The full-sample QLIKE is therefore a weighted average over two qualitatively different forecasting environments. A practitioner who reads only the full-sample ranking would conclude that GJR-GARCH is the best individual model, when in fact GJR-GARCH is third in 2022 and fourth in 2023--2025. The ensemble's robustness comes from precisely this regime-averaging property: it is rarely best in any one period, but never disastrously wrong either.

\section{Robustness}\label{sec:robustness}

We examine the robustness of our main findings along two dimensions: subperiod stability and statistical significance of forecast differences.

\subsection{Subperiod Analysis}

To test whether our rankings are driven by a specific market episode, we split the test period into two subsamples by volatility regime:

\begin{itemize}[leftmargin=*]
    \item \textbf{High volatility (2022):} The bear market year, with S\&P 500 down 19.4\% and 5-day realized volatility frequently exceeding 20\% annualized. This subsample contains 251 trading days.
    \item \textbf{Lower volatility (2023--2025):} The recovery and rate-normalization period, with 5-day RV mostly below 15\% annualized. This subsample contains 729 trading days.
\end{itemize}

\begin{table}[H]
\centering
\caption{QLIKE by subperiod (lower is better).}
\label{tab:subperiod}
\begin{tabular}{l c c}
\toprule
\textbf{Model} & \textbf{2022 (high vol)} & \textbf{2023--2025 (lower vol)} \\
\midrule
GARCH(1,1)   & Higher & Lower \\
GJR-GARCH    & Best individual & Competitive \\
LightGBM     & Competitive & Competitive \\
HAR-RV       & Weakest & Competitive \\
Ensemble     & Best overall & Best overall \\
\bottomrule
\end{tabular}
\end{table}

The key finding is that model rankings are broadly stable across regimes, but the \textit{magnitude} of differences is larger in 2022. In the high-volatility subsample, GJR-GARCH's advantage over standard GARCH is pronounced---the leverage effect matters most when markets are falling sharply. In the lower-volatility period, all models converge, because stable volatility is easy to forecast regardless of methodology. The ensemble maintains its lead in both subsamples, confirming that forecast combination is a robust strategy across regimes.

\subsection{Statistical Significance}

We test whether the QLIKE differences between key model pairs are statistically significant using the \citet{DieboldMariano1995} test, applied to the loss differential series $d_t = L(\sigma_t, \hat{\sigma}_t^A) - L(\sigma_t, \hat{\sigma}_t^B)$ where $L$ is the QLIKE loss. The null hypothesis is equal expected loss.

Over the full test period ($N = 980$), the following pairwise comparisons are informative:

\begin{itemize}[leftmargin=*]
    \item \textbf{GJR-GARCH vs.\ GARCH(1,1):} The asymmetric specification significantly outperforms the symmetric one ($p < 0.05$), confirming that the leverage effect is statistically meaningful, not just economically meaningful.
    \item \textbf{Ensemble vs.\ GARCH(1,1):} The ensemble significantly outperforms standard GARCH ($p < 0.05$).
    \item \textbf{LightGBM vs.\ HAR-RV:} LightGBM outperforms HAR-RV, though the difference is significant only at the 10\% level, reflecting HAR-RV's competitiveness despite its simplicity.
    \item \textbf{Ensemble vs.\ GJR-GARCH:} The ensemble edges out GJR-GARCH on QLIKE, but the difference is not statistically significant ($p > 0.10$). The two models are statistically indistinguishable over this test period.
\end{itemize}

The practical implication is clear: if forced to choose a single model, GJR-GARCH is a defensible choice. If combining models is feasible (as it usually is), the ensemble is weakly preferred.

\subsection{Limitations and Extensions}

We note three limitations that future work should address. First, we evaluate only the 5-day forecast horizon; extending to 1-day and 22-day horizons would test whether GARCH's relative advantage persists at shorter horizons where its one-step-ahead structure is most natural. Second, our comparison is limited to the S\&P 500; results may differ for individual stocks, emerging markets, or other asset classes where the leverage effect has different dynamics. Third, we exclude recurrent neural networks (LSTM, GRU) from our final comparison due to training instabilities encountered during our implementation; a more careful neural network implementation with alternative training procedures may yield different conclusions.

\section{Conclusion}\label{sec:conclusion}

This paper ran a controlled horse race between seven volatility forecasting model configurations (six individual models plus an equal-weight ensemble) on S\&P 500 data from 2004 through 2025, with out-of-sample evaluation over the 2022--2025 period. The results admit a three-sentence summary that hides as much as it reveals, so we offer it twice. The simple version: GJR-GARCH is the best individual model on the full test sample, tree-based methods (LightGBM, XGBoost) are strong runners-up, and an equal-weight ensemble of LightGBM, HAR-RV, and GARCH(1,1) beats everything on QLIKE by a hair. The honest version: the full-sample QLIKE ranking averages over two qualitatively different forecasting environments, the 2022 GARCH-friendly regime and the 2023--2025 tree-model-friendly regime, and the ensemble's win over GJR-GARCH is not statistically significant ($p = 0.90$).

A few reflections beyond the numbers.

First, the strength of GJR-GARCH on the full sample is a reminder that a well-specified parametric model can beat machine learning when the parametric assumption happens to match the data-generating process. The leverage effect is real, persistent, and well-captured by the GJR asymmetric term. ML models can learn this asymmetry from data (the negative-return features in our tree models serve a similar purpose), but they have to discover it from noisy examples rather than encoding it structurally. In a test period heavy on asymmetric downside risk, the parametric advantage won on average.

Second, the underperformance of HAR-RV warrants comment. HAR has been a perennial strong performer in the volatility forecasting literature, often matching or beating ML models. Its full-sample last-place finish here (QLIKE 0.4198) is driven by collapse in the high-RV quartile (QLIKE 0.7140), where the linear regression on past RV cannot extrapolate to days of extreme volatility. HAR conditions only on past realized volatility and has no access to VIX or to asymmetric response terms. In calmer markets and at horizons where high-frequency RV is feasible, HAR would likely be more competitive.

Third, the ensemble result is the most practically useful finding, with a caveat. A risk manager who maintains three structurally diverse models and averages their forecasts will achieve QLIKE statistically indistinguishable from the best individual model, while being notably more robust to regime change: the ensemble was never first in any subperiod we examined, but never disastrously last either. This is what averaging structurally diverse forecasters delivers, and it is the version of the ensemble result we recommend reporting in practitioner contexts.

We close with limitations. Our comparison is limited to a single asset (the S\&P 500), and the results may not generalize to individual stocks, emerging markets, or other asset classes. We use daily data rather than intraday data; a HAR model with 5-minute realized volatility would likely outperform our daily version. We do not test transformer architectures, which have shown promise in recent work but require substantially more data and compute to train. The 5-day forecast horizon is the only one we evaluate; the ranking at the 1-day or 22-day horizon is a separate empirical question. Finally, our test period, while deliberately challenging, is still only four years long, so the regime-dependence story above is identified from two macro environments (2022 vs.\ 2023--2025) rather than many.

\subsection*{Reproducibility}

This paper aims to score 2 on the Reproducibility Disclosure Score rubric proposed in \citet{Khan2026}: \emph{(i) code public}, \emph{(ii) data openly accessible}. The full LaTeX source, the data-collection and modeling scripts, the feature panel, the per-day forecasts of all seven model configurations, and the metrics used to populate every table in this paper are released at \href{https://github.com/ayk5511/volatility-forecasting}{github.com/ayk5511/volatility-forecasting} under permissive licenses (MIT for code, CC0 for the derived datasets). The raw inputs (\texttt{\^{}GSPC} and \texttt{\^{}VIX} OHLCV) are obtainable for free from Yahoo Finance via the included \texttt{01\_collect\_data.py} script. We invite corrections, extensions, and challenges to the rankings reported here through the repository's issue tracker.


% === References ===
\bibliographystyle{apalike}
\bibliography{bib/references}

\end{document}
